\chapter{Benefícios}

\section{Benefícios Diretos para os Participantes}

\begin{itemize}
    \item Acesso a uma ferramenta educacional gamificada que facilita a preparação para a prova de matemática do ENEM;
    \item Desenvolvimento de habilidades cognitivas (decomposição, abstração, reconhecimento de padrões e algoritmos) por meio da resolução de questões reais;
    \item Experiência de aprendizado ativa e psicologicamente estimulante, substituindo a postura passiva e a frustração comum no estudo de matemática.
\end{itemize}

\section{Benefícios Indiretos para a Sociedade}

\begin{itemize}
    \item Contribuição para o campo de Educação Computacional e Design-Based Research, com documentação aberta do processo de desenvolvimento e validação;
    \item Geração de dados educacionais que podem alimentar análises de mineração de dados educacionais em trabalhos futuros;
    \item Modelo replicável de integração entre PC e gamificação, alinhado às diretrizes da BNCC e da BNCC Computação;
    \item Potencial de mitigação do baixo desempenho em matemática no ENEM, impactando positivamente o acesso ao ensino superior.
\end{itemize}

\section{Forma de Aplicação dos Benefícios}

\begin{itemize}
    \item Disponibilização da plataforma Cognitivo para uso dos participantes durante a pesquisa;
    \item Retorno dos resultados da pesquisa aos participantes e à instituição, por meio de relatório e/ou apresentação acadêmica.
\end{itemize}